%!TEX root = ../template.tex
%%%%%%%%%%%%%%%%%%%%%%%%%%%%%%%%%%%%%%%%%%%%%%%%%%%%%%%%%%%%%%%%%%%%
%% chapter6.tex
%% NOVA thesis document file
%%
%% Chapter with lots of dummy text
%%%%%%%%%%%%%%%%%%%%%%%%%%%%%%%%%%%%%%%%%%%%%%%%%%%%%%%%%%%%%%%%%%%%

\typeout{NT FILE chapter6.tex}%

\chapter{Translation}
\label{cha:Translation}

%%maybe rewrite, probably too big for the first part

For the following chapter it is referred the creation from base the full translation of \cml with the new \gospel annotations that were
previously not present in the code. Unlike extraction, which derives executable code from a verified program by eliminating its logical 
content, translation operates in the opposite direction, taking an existing implementation together with its specifications and expressing 
both in \whyml so that verification can be carried out within the Why3 platform. The purpose of this process is not to replace extraction 
but to serve as a basis for comparison: by independently constructing a \whyml representation of the same program, it becomes possible to 
evaluate the two approaches against each other in terms of the verification effort required, the expressiveness of the resulting specifications, 
and the degree to which each method preserves the intent of the original code. The following section details the translation process and the 
decisions made when mapping \cml and \gospel constructs into their \whyml counterparts.

\section{Lexer}
\label{sec:Lexer}

Before any syntactic or semantic analysis can take place, the raw source text must be broken down into meaningful units. This is the responsibility 
of the lexer, which scans the input sequentially and produces a stream of tokens that the parser can then process. Every raw source text that it is
receiving regarding \cml it will deconstruct it into tokens for the parser, on the other hand the raw text received inside the delimited \gospel
specifications, in between the \inlinecode{(*@} and \inlinecode{*)}, will just return the raw text into parser so it can be deconstructed in another
phase later.

\section{Parser}
\label{sec:Parser}

The parser processes the token stream produced by the lexer and constructs an abstract syntax tree representing the input program. Each recognised 
is mapped to a corresponding AST node, producing a structured representation that the subsequent translation phase can operate on directly. 
Operator precedence and associativity are resolved at this stage in accordance with \cml's semantics, ensuring that the structure of expressions 
is faithfully preserved prior to translation into \whyml. Lastly, \gospel annotations are still preserved as raw strings in the abstract syntax tree, 
deferring their interpretation to a later stage of the translation process. The resulting AST serves as the central intermediate representation 
from which the translation into \whyml is driven.

\section{Abstract Syntax Tree(AST)}
\label{sec:AST}

The AST serves as the first structured internal representation of the input program, produced directly from the parsing phase. At this stage, the 
AST captures the full syntactic content of the \cml source file, each represented as a typed node, that reflects the grammatical structure of the input. 
\gospel annotations, having been recognised by the lexer and passed through the parser as raw strings, are preserved at this stage without further 
interpretation, attached to their corresponding top-level positions in the tree.

\section{\gospel Abstract Syntax Tree(GAST)}
\label{sec:GAST}

The GAST represents the final structured intermediate representation prior to translation into \whyml. It extends the AST by replacing raw \gospel 
annotation strings with fully parsed \gospel nodes, producing a unified representation in which \cml program content and \gospel specifications 
coexist in a structured and typed form.

\section{AST To GAST}
\label{sec:AST_to_GAST}

The transformation from AST to GAST is responsible for promoting the raw \gospel annotation strings preserved in the AST into structured \gospel AST nodes, 
producing a unified representation in which each \cml definition is paired with its corresponding specification. This phase is driven by the \gospel parser, 
which is invoked selectively depending on the nature of each annotation encountered.

A central distinction in this transformation is between \textit{floating} and \textit{inline} \gospel annotations. Floating annotations are top-level logical 
declarations that stand independently of any function definition, these include predicates, functions, axioms, and lemmas, which are parsed and mapped directly 
to their corresponding GAST nodes. Inline annotations, by contrast, are specifications attached to function definitions, such as pre-conditions and post-conditions,
that are associated with the definition that follows or precedes them in the source file. To handle this, consecutive inline annotation strings are collected and 
merged before being passed to the \gospel parser as a single unified specification.

A further consideration arises from the fact that \gospel annotations may appear either before or after their associated function definition in the source file. 
The transformation accounts for both orderings, collecting specification fragments from either side of the definition and combining them into a single coherent 
specification prior to parsing.

\section{Printer}
\label{sec:Printer}

The final phase of the translation is the printer, where traversing the GAST and printing the correct \whyml code. While much of the translation is a 
straightforward structural mapping from \cml constructs to their \whyml equivalents, several design decisions merit discussion from a verification standpoint.

Another decision to note is the automatic detection of recursive functions. Rather than requiring the user to explicitly mark recursive definitions, the 
printer analyses the function body to determine whether it contains a self-referential call, emitting the \inlinecode{rec} keyword accordingly.
This analysis correctly accounts for shadowed identifiers, ensuring that a let binding or anonymous function argument sharing the function's name 
does not incorrectly trigger a recursive annotation.

Finally, the handling of \inlinecode{Absurd} in the printer is consistent with the decision taken in the extraction driver. Whereas extraction maps 
\inlinecode{absurd} to \inlinecode{raise Absurd}, the printer maps \inlinecode{raise Absurd} back to \inlinecode{absurd} in \whyml, restoring the original 
construct's semantics and ensuring that unreachable branches are correctly represented in the verification context.